% Part: lambda-calculus
% Chapter: syntax
% Section: de-bruijn

\documentclass[../../../include/open-logic-section]{subfiles}

\begin{document}

\olfileid{lam}{syn}{deb}

\olsection{نمایهٔ دوبروین}

هم‌ارزی~$\alpha$ بسیار طبیعی است، زیرا ترم‌های هم‌ارز به~$\alpha$
«معنای یکسانی دارند». در واقع می‌توان نحوی برای ترم‌های لامبدا ارائه کرد
که میان ترم‌هایی که با تبدیل~$\alpha$ به یکدیگر تبدیل می‌شوند تمایز
نمی‌گذارد. شناخته‌شده‌ترین شیوه، متغیرها را با \emph{نمایهٔ دوبروین}
آن‌ها جایگزین می‌کند.

وقتی~$\lambd[x][M]$ را می‌نویسیم، به‌صراحت بیان می‌کنیم که~$x$ پارامتر
تابع است، تا در~$M$ با استفاده از~$x$ به این پارامتر ارجاع دهیم. اما در
نمایهٔ دوبروین، پارامترها نام ندارند و ارجاع به آن‌ها در بدنهٔ تابع با
عددی نشان داده می‌شود که شمار سطوح انتزاع میان ارجاع و پارامتر را بیان
می‌کند. برای مثال، ترم~$\lambd[x][\lambd[y][y x]]$ را در نظر بگیرید:
انتزاع بیرونی متغیر~$x$ را مقید می‌کند؛ انتزاع درونی متغیر~$y$ را مقید
می‌کند؛ و زیرترم~$y x$ در دامنهٔ انتزاع درونی قرار دارد: میان~$y$ و
انتزاع‌کنندهٔ آن~$\lambd[y]$ هیچ انتزاعی نیست، اما میان~$x$ و
انتزاع‌کنندهٔ آن~$\lambd[x]$ یک انتزاع هست. پس برای~$y x$ می‌نویسیم
$0\, 1$، و برای کل ترم می‌نویسیم~$\lambd[][\lambd[][01]]$.

\begin{defn}
  ترم‌های دوبروین به‌صورت استقرایی چنین تعریف می‌شوند:
  \begin{enumerate}
  \item $n$، که در آن~$n$ هر عدد طبیعی است.
  \item $PQ$، که در آن~$P$ و~$Q$ هر دو ترم دوبروین‌اند.
  \item $\lambd[][N]$، که در آن~$N$ یک ترم دوبروین است.
  \end{enumerate}
\end{defn}

ترجمه‌ای صوری از ترم‌های معمول لامبدا به ترم‌های نمایه‌گذاری‌شدهٔ
دوبروین چنین است:
\begin{defn}
  \begin{align*}
    F_\Gamma(x) &= \Gamma(x) \\
    F_\Gamma(PQ) &= F_\Gamma(P)F_\Gamma(Q) \\
    F_\Gamma(\lambd[x][N]) &= \lambd[][F_{x,\Gamma}(N)]
  \end{align*}
  که در آن~$\Gamma$ فهرستی از متغیرهاست که از صفر نمایه‌گذاری شده، و
  $\Gamma(x)$ جایگاه متغیر~$x$ را در~$\Gamma$ نشان می‌دهد. برای مثال،
  اگر~$\Gamma$ برابر با~$x,y,z$ باشد، $\Gamma(x)$ برابر با~$0$ و
  $\Gamma(z)$ برابر با~$2$ است.
  
  نماد~$x,\Gamma$ فهرست حاصل از افزودن~$x$ به ابتدای~$\Gamma$ را نشان
  می‌دهد؛ برای نمونه، با ادامهٔ~$\Gamma$ مثال پیشین، $w,\Gamma$ برابر
  با~$w,x,y,z$ است.
\end{defn}

بازیابی یک ترم معیار لامبدا از یک ترم دوبروین چنین انجام می‌شود:

\begin{defn}
  \begin{align*}
    G_\Gamma(n) &= \Gamma[n] \\
    G_\Gamma(PQ) &= G_\Gamma(P) G_\Gamma(Q) \\
    G_\Gamma(\lambd[][N]) &= \lambd[x][G_{x,\Gamma}(N)]
  \end{align*}
  که در آن~$\Gamma$ باز فهرستی از متغیرهاست که از صفر نمایه‌گذاری شده،
  و~$\Gamma[n]$ متغیر جایگاه~$n$ را نشان می‌دهد. برای مثال، اگر~$\Gamma$
  برابر با~$x,y,z$ باشد، آنگاه~$\Gamma[1]$ برابر با~$y$ است.

  متغیر~$x$ در معادلهٔ آخر هر متغیری برگزیده می‌شود که در~$\Gamma$
  نباشد.
\end{defn}

در اینجا چند نتیجه را بی‌اثبات می‌آوریم:

\begin{prop}
  اگر~$M \aconvone M'$ و~$\Gamma$ هر فهرستی شامل~$\FV{M}$ باشد، آنگاه
  $F_\Gamma(M) \eqs F_\Gamma(M')$.
\end{prop}

\end{document}
